\chapter{Introdução}
\paragraph{Modelo da Máquina}
O modelo da máquina CC é descrito pelas seguintes equações de estado:
\begin{subequations}

$
X=
\begin{bmatrix}
I_{a} \\
W_{r}
\end{bmatrix}
$

~

$
U=
\begin{bmatrix}
V_{a}\\
T_{m}
\end{bmatrix}
$

~

$
\begin{bmatrix}
\dot{I}_{a} \\
\dot{W}_{r}
\end{bmatrix}
=
\begin{bmatrix}
-\frac{R_{a}}{L_{a}}-\frac{K_{e}\lambda_{e}}{L_{a}} \\
\frac{K_{e}\lambda_{e}}{J_{m}}-\frac{F_{m}}{J_{m}}
\end{bmatrix}
\begin{bmatrix}
I_{a}\\
W_{r}
\end{bmatrix}
+
\begin{bmatrix}
1/L_{a} & 0 \\
0 & -1/J_{m}
\end{bmatrix}
\begin{bmatrix}
V_{a}\\
T_{m}
\end{bmatrix}
$
\end{subequations}

Escolhemos controlar a máquina a partir da tensão de armadura $V_{a}$ e projetamos um controlador tendo em mente o cancelamento de polos. Usamos os seguintes parâmetros para a máquina:

$
K_{e}=476.16 ; R_{a}=0.08\Omega ; L_{a}=0.0001H ; I_{e}=2A ;
J_{m}=0.001Kg.m^2 ; F_{m}=0.0001N.m.s/rad
$


\paragraph{Modelo do Controlador com parâmetros conhecidos}

Assumindo que os parâmetros da planta são conhecidos calculamos os valores de $K_{p}, K_{i}, K_{d}$, assumindo $T_{d}$ da parte derivativa como sendo 0.005. Desta forma o controlador projetado para o caso em que os parâmetros da planta são conhecidos tem valores:


$ K_{p}=-0.2339 ; K_{i}=47.6164; K_{d}=1.1747$, e portanto em malha fechada:
~

$ G(s)=\big(\frac{9999.444}{s^2 + 200s + 9999.444}) $
